\documentclass{jsarticle}
\usepackage{amsmath,amssymb,amsthm,bm,physics,arydshln}
\newtheorem{question}{問}
\newtheorem*{answer*}{解答}
\begin{document}

\begin{question}
次の行列のJordan標準形を求めよ.
\end{question}

(1)
$\left(\begin{array}{cc}
1 & 4 \\
-1 & -3
\end{array}\right)$ \ \ \ \ \ \ \ \ \ \ 
(2)
\( \left(\begin{array}{ccccccc}0 & 1 & 0 & -1 & -1 & 0 & 1 \\ 0 & 0 & 1 & 1 & 1 & 0 & -1 \\ 0 & 0 & 0 & 0 & -1 & 0 & 1 \\ 0 & 0 & 0 & 0 & 1 & 0 & -1 \\ 0 & 0 & 0 & 0 & 0 & 1 & 1 \\ 0 & 0 & 0 & 0 & 0 & 0 & 0 \\ 0 & 0 & 0 & 0 & 0 & 0 & 0\end{array}\right) \) \ \ \ \ \ \ \ \ \ \ 
(3)
$
\left(\begin{array}{ccccc}
1 & 1 & \cdots & \cdots & 1 \\
& 0 & & & \vdots \\
& & \ddots & & \vdots \\
& & & 0 & 1 \\
& & & & 1
\end{array}\right)
$




\begin{answer*}
\end{answer*}
(1)

(解法 1)
\[
|t E-A|=\left|\begin{array}{cc}
t-1 & -4 \\
1 & t+3
\end{array}\right|=(t-1)(t+3)+4=t^2+2 t+1=(t+1)^2 \text { なので$A$の固有値は$-1$} .
\]
\( \begin{aligned} & P:=\left((A+E)\left(\begin{array}{ll}1 \\ 0\end{array}\right) \ \   \left(\begin{array}{l}1 \\ 0\end{array}\right)\right)=\left(\begin{array}{cc}2 & 1 \\ -1 & 0\end{array}\right) とおくと, \\ & P^{-1} A P=\left(\begin{array}{cc}0 & -1 \\ 1 & 2\end{array}\right)\left(\begin{array}{cc}1 & 4 \\ -1 & -3\end{array}\right)\left(\begin{array}{cc}2 & 1 \\ -1 & 0\end{array}\right)=\left(\begin{array}{cc}1 & 3 \\ -1 & -2\end{array}\right)\left(\begin{array}{cc}2 & 1 \\ -1 & 0\end{array}\right)=\left(\begin{array}{cc}-1 & 1 \\ 0 & -1\end{array}\right)\end{aligned} \)

\bigskip
\bigskip

（解法2）

固有値は$-1$なので，$A+E$ \( =\left(\begin{array}{cc}2 & 4 \\ -1 & -2\end{array}\right) \rightarrow\left(\begin{array}{ll}1 & 2 \\ 0 & 0\end{array}\right) \) より, 
$(A+E) \bm{x_1}=\bm{0}$ を解くと, 

\(\quad \bm{x_1}=\left(\begin{array}{c}-2 c_1 \\ c_1\end{array}\right)=c_1\left(\begin{array}{c}-2 \\ 1\end{array}\right). \quad \therefore W(-1 ; A)=\left\langle\left(\begin{array}{c}-2 \\ 1\end{array}\right)\right\rangle \).

\( (A+E) \bm{x_2}=\bm{x_1} を解くと, \quad \bm{x_2}=\left(\begin{array}{c}-2 c_2-1 \\ c_2\end{array}\right). \quad c_2=0 とすると,  \quad x_2=\left(\begin{array}{c}-1 \\ 0\end{array}\right) \).

$P:=\left(\begin{array}{cc}-2 & -1 \\ 1 & 0\end{array}\right)$ とおくことで


\( \begin{aligned} P^{-1} A P=\left(\begin{array}{cc}0 & 1 \\ -1 & -2\end{array}\right)\left(\begin{array}{cc}1 & 4 \\ -1 & -3\end{array}\right)\left(\begin{array}{rr}-2 & -1 \\ 1 & 0\end{array}\right) & =\left(\begin{array}{cc}-1 & -3 \\ 1 & 2\end{array}\right)\left(\begin{array}{cc}-2 & -1 \\ 1 & 0\end{array}\right) \\ & =\left(\begin{array}{cc}-1 & 1 \\ 0 & -1\end{array}\right)\end{aligned} \)

\bigskip
\bigskip

(2)

(解法１)

固有多項式は$t^7$ で$A^2 \neq O , \ A^3 = O$ なので, 最小多項式は$t^3$. 

よって, 固有値$0$のJordan細胞の最大次数は$3$. 

また, 
$
\left(\begin{array}{ccccccc:c}
0 & 1 & 0 & -1 & -1 & 0 & 1 & b_1 \\
0 & 0 & 1 & 1 & 1 & 0 & -1 & b_2 \\
0 & 0 & 0 & 0 & -1 & 0 & 1 & b_3 \\
0 & 0 & 0 & 0 & 1 & 0 & -1 & b_4 \\
0 & 0 & 0 & 0 & 0 & 1 & 1 & b_5 \\
0 & 0 & 0 & 0 & 0 & 0 & 0 & b_6 \\
0 & 0 & 0 & 0 & 0 & 0 & 0 & b_7
\end{array}\right) \rightarrow\left(\begin{array}{ccccccc:c}
0 & 1 & 1 & 0 & 0 & 0 & 0 & b_1+b_2 \\
0 & 0 & 1 & 1 & 0 & 0 & 0 & b_2+b_3 \\
0 & 0 & 0 & 0 & 1 & 0 & -1 & -b_3 \\
0 & 0 & 0 & 0 & 0 & 1 & 1 & b_5 \\
0 & 0 & 0 & 0 & 0 & 0 & 0 & b_3+b_4 \\
0 & 0 & 0 & 0 & 0 & 0 & 0 & b_6 \\
0 & 0 & 0 & 0 & 0 & 0 & 0 & b_7
\end{array}\right)
$
なので, 

\bigskip

dim(W(0;$A$)) = $3$. \ $\therefore$ 固有値$0$のJordan細胞の個数は$3$.

求めるJordan標準形は
$J(0,3) \oplus J(0,3) \oplus J(0,1) \quad または \quad J(0,3) \oplus J(0,2) \oplus J(0,2)$.

$A\bm{x} = \bm{b}$ が解をもつための条件は, $b_3+b_4 = 0, \ b_6=b_7=0$ ────(*)

$A\bm{x} = \bm{0}$ を解くことにより, $W(0;A) = \begin{pmatrix} a_1 \\ 0 \\ a_2 \\ -a_2 \\ a_3 \\ -a_3 \\ a_3 \end{pmatrix} = a_1 \begin{pmatrix} 1 \\ 0 \\ 0 \\ 0 \\ 0 \\ 0 \\ 0 \end{pmatrix} + a_2 \begin{pmatrix} 0 \\ 1 \\ -1 \\ 1 \\ 0 \\ 0 \\ 0 \end{pmatrix} + a_3 \begin{pmatrix} 0 \\ 0 \\ 0 \\ 0 \\ 1 \\ -1 \\ 1 \end{pmatrix}$

○○は(*)を満たさないので, $A$のJordan標準形は$1$次のJordan細胞を持つことがわかる.

よって, Jordan標準形は$\underline{J(0,3) \oplus J(0,3) \oplus J(0,1)}$.

\bigskip
\bigskip

$A\bm{x} = \bm{0}$ を解くことにより, $W(0;A) = \begin{pmatrix} a_1 \\ 0 \\ a_2 \\ -a_2 \\ a_3 \\ -a_3 \\ a_3 \end{pmatrix} = a_1 \begin{pmatrix} 1 \\ 0 \\ 0 \\ 0 \\ 0 \\ 0 \\ 0 \end{pmatrix} + a_2 \begin{pmatrix} 0 \\ 1 \\ -1 \\ 1 \\ 0 \\ 0 \\ 0 \end{pmatrix} + a_3 \begin{pmatrix} 0 \\ 0 \\ 0 \\ 0 \\ 1 \\ -1 \\ 1 \end{pmatrix}$

$A\bm{x} = \begin{pmatrix} 1 \\ 0 \\ 0 \\ 0 \\ 0 \\ 0 \\ 0 \end{pmatrix}$ を解くことにより, 
$\bm{x} = \begin{pmatrix} 0 \\ 0 \\ 1 \\ -1 \\ 0 \\ 0 \\ 0 \end{pmatrix} + a_4\begin{pmatrix} 1 \\ 0 \\ 0 \\ 0 \\ 0 \\ 0 \\ 0 \end{pmatrix} + a_5\begin{pmatrix} 0 \\ 1 \\ -1 \\ 1 \\ 0 \\ 0 \\ 0 \end{pmatrix} + a_6\begin{pmatrix} 0 \\ 0 \\ 0 \\ 0 \\ 1 \\ -1 \\ 1 \end{pmatrix}$

$A\bm{x} = \begin{pmatrix} 0 \\ 1 \\ 0 \\ 0 \\ 0 \\ 0 \\ 0 \end{pmatrix}$ を解くことにより, 
$\bm{x} = \begin{pmatrix} 0 \\ 0 \\ 1 \\ 0 \\ 0 \\ 0 \\ 0 \end{pmatrix} + a_7\begin{pmatrix} 1 \\ 0 \\ 0 \\ 0 \\ 0 \\ 0 \\ 0\end{pmatrix} + a_8\begin{pmatrix} 0 \\ 1 \\ -1 \\ 1 \\ 0 \\ 0 \\ 0 \end{pmatrix} + a_9\begin{pmatrix} 0 \\ 0 \\ 0 \\ 0 \\ 1 \\ -1 \\ 1 \end{pmatrix}$

$A\bm{x} = \begin{pmatrix} 0 \\ 1 \\ -1 \\ 1 \\ 0 \\ 0 \\ 0 \end{pmatrix}$ を解くことにより, 
$\bm{x} = \begin{pmatrix} 0 \\ 0 \\ 1 \\ -1 \\ 0 \\ 1 \\ -1 \end{pmatrix} + a_{10} \begin{pmatrix} 1 \\ 0 \\ 0 \\ 0 \\ 0 \\ 0 \\ 0 \end{pmatrix} + a_{11} \begin{pmatrix} 0 \\ 1 \\ -1 \\ 1 \\ 0 \\ 0 \\ 0 \end{pmatrix} + a_{12} \begin{pmatrix} 0 \\ 0 \\ 0 \\ 0 \\ 1 \\ -1 \\ 1 \end{pmatrix}$

$A\bm{x} = \begin{pmatrix} 0 \\ 0 \\ 1 \\ -1 \\ 1 \\ 0 \\ 0 \end{pmatrix}$ を解くことにより,
$\bm{x} = \begin{pmatrix} 0 \\ 0 \\ 0 \\ 1 \\ 0 \\ 0 \\ 1 \end{pmatrix} + a_{13} \begin{pmatrix} 1 \\ 0 \\ 0 \\ 0 \\ 0 \\ 0 \\ 0 \end{pmatrix} + a_{14} \begin{pmatrix}  0 \\ 1 \\ -1 \\ 1 \\ 0 \\ 0 \\ 0 \end{pmatrix} + a_{15} \begin{pmatrix} 0 \\ 0 \\ 0 \\ 0 \\ 1 \\ -1 \\ 1 \end{pmatrix}$

$P:= \begin{pmatrix} 1 & 0 & 0 & 0 & 0 & 0 & 0 \\ 0 & 1 & 0 & 1 & 0 & 0 & 0 \\ 0 & 0 & 1 & -1 & 1 & 0 & 0 \\ 0 & 0 & 0 & 1 & -1 & 1 & 0 \\ 0 & 0 & 0 & 0 & 1 & 0 & 1 \\ 0 & 0 & 0 & 0 & 0 & 0 & -1 \\ 0 & 0 & 0 & 0 & 0 & 1 & 1 \end{pmatrix}$ とおくと, 
$P^{-1} = \begin{pmatrix} 1 & 0 & 0 & 0 & 0 & 0 & 0 \\ 0 & 1 & 0 & -1 & -1 & 0 & 1 \\ 0 & 0 & 1 & 1 & 0 & -1 & -1 \\ 0 & 0 & 0 & 1 & 1 & 0 & -1 \\ 0 & 0 & 0 & 0 & 1 & 1 & 0 \\ 0 & 0 & 0 & 0 & 0 & 1 & 1 \\ 0 & 0 & 0 & 0 & 0 & -1 & 0 \end{pmatrix}$ であり, 

$P^{-1}AP = \begin{pmatrix} 0 & 1 & 0 & -1 & -1 & 0 & 1 \\ 0 & 0 & 1 & 1 & 0 & -1 & -1 \\ 0 & 0 & 0 & 0 & 0 & 0 & 0 \\ 0 & 0 & 0 & 0 & 1 & 1 & 0 \\ 0 & 0 & 0 & 0 & 0 & 1 & 1 \\0 & 0 & 0 & 0 & 0 & 0 & 0 \\ 0 & 0 & 0 & 0 & 0 & 0 & 0 \end{pmatrix} \begin{pmatrix} 1 & 0 & 0 & 0 & 0 & 0 & 0 \\ 0 & 1 & 0 & 1 & 0 & 0 & 0 \\ 0 & 0 & 1 & -1 & 1 & 0 & 0 \\ 0 & 0 & 0 & 1 & -1 & 1 & 0 \\ 0 & 0 & 0 & 0 & 1 & 0 & 1 \\ 0 & 0 & 0 & 0 & 0 & 0 & -1 \\ 0 & 0 & 0 & 0 & 0 & 1 & 1 \end{pmatrix} = J(0,3) \oplus J(0,3) \oplus J(0,1)$

(解法2)

$xE - A \rightarrow \left(\begin{array}{ccccccc}1 & 0 & 0 & 0 & 0 & 0 & 0 \\ 0 & 1 & 0 & 0 & 0 & 0 & 0 \\ 0 & 0 & 1 & 0 & 0 & 0 & 0 \\ 0 & 0 & 0 & 1 & 0 & 0 & 0 \\ 0 & 0 & 0 & 0 & x & 0 & 0 \\ 0 & 0 & 0 & 0 & 0 & x^3 & 0 \\ 0 & 0 & 0 & 0 & 0 & 0 & x^3\end{array}\right) \rightarrow \left(\begin{array}{c:ccc:ccc}x & 0 & 0 & 0 & 0 & 0 & 0 \\ \hdashline 0 & 1 & 0 & 0 & 0 & 0 & 0 \\ 0 & 0 & 1 & 0 & 0 & 0 & 0 \\ 0 & 0 & 0 & x^3 & 0 & 0 & 0 \\ \hdashline 0 & 0 & 0 & 0 & 1 & 0 & 0 \\ 0 & 0 & 0 & 0 & 0 & 1 & 0 \\ 0 & 0 & 0 & 0 & 0 & 0 & x^3\end{array}\right)$

$\hspace{20zw} \sim 
\left(\begin{array}{c:ccc:ccc}x & 0 & 0 & 0 & 0 & 0 & 0 \\ \hdashline 0 & x-1 & 0 & 0 & 0 & 0 & 0 \\ 0 & 0 & x-1 & 0 & 0 & 0 & 0 \\ 0 & 0 & 0 & x & 0 & 0 & 0 \\ \hdashline 0 & 0 & 0 & 0 & x-1 & 0 & 0 \\ 0 & 0 & 0 & 0 & 0 & x-1 & 0 \\ 0 & 0 & 0 & 0 & 0 & 0 & x\end{array}\right)$

$\hspace{20zw} = 
xE_7 - \left(\begin{array}{c:ccc:ccc}0 & 0 & 0 & 0 & 0 & 0 & 0 \\ \hdashline 0 & 0 & 1 & 0 & 0 & 0 & 0 \\ 0 & 0 & 0 & 1 & 0 & 0 & 0 \\ 0 & 0 & 0 & 0 & 0 & 0 & 0 \\ \hdashline 0 & 0 & 0 & 0 & 0 & 1 & 0 \\ 0 & 0 & 0 & 0 & 0 & 0 & 1 \\ 0 & 0 & 0 & 0 & 0 & 0 & 0\end{array}\right)$

なので, Jordan標準形は $J(0,3) \oplus J(0,3) \oplus J(0,1)$.

\bigskip
\bigskip

\begin{question}
次の行列をある正則行列によって対角化せよ.
\end{question}

(1)
$\left(\begin{array}{cccc}
1 & \cdots & \cdots & 1 \\
\vdots & & & \\
\vdots & & \text{\huge{O}} & \\
1 & & & \end{array} \right)$

\begin{answer*}
\end{answer*}
固有多項式は, 
\begin{align*}
D_n(t) =
\begin{vmatrix}
t-1 & -1 & -1 & \cdots & -1 \\
-1 & t & 0 & \cdots & 0 \\
-1 & 0 & t & & \\
\vdots & \vdots & & \ddots \\
-1 & 0 & & & t
\end{vmatrix}
&= \begin{vmatrix}
t & -1 & 0 & \cdots & 0 \\
-1 & t-1 & -1 & \cdots & -1 \\
0 & -1 & t & & \\
\vdots & \vdots & & \ddots & \\
0 & -1 & & & t
\end{vmatrix} \\
&= t\begin{vmatrix} 
t-1 & -1 & \cdots & -1 \\
-1 & t & &  \\
\vdots & & \ddots & \\
-1 & & & t
\end{vmatrix}
+ \begin{vmatrix}
-1 & 0 & \cdots & 0 \\
-1 & t & & \\
\vdots & & \ddots & \\
-1 & & & t
\end{vmatrix} \ (余因子展開) \\
&= tD_{n-1}(t) - t^{n-2} \ (n \ge 3)
\end{align*}

(i) $t = 0$ のとき, $D_n(0) = 0$.

(ii) $t \neq 0$ のとき, 

$\displaystyle \frac{D_n(x)}{x^n} = \displaystyle \frac{D_{n-1}(x)}{x^{n-1}} - \displaystyle \frac{1}{x^2}$ なので, 漸化式を解くと
$\displaystyle \frac{D_n(x)}{x^n} = \displaystyle \frac{x^2 - x - n + 1}{x^2}$. 

(i), (ii) より$D_n(x) = x^{n-2}\left(x - \displaystyle \frac{1 + \sqrt{4n-3}}{2}\right) \left(x - \displaystyle \frac{1 - \sqrt{4n-3}}{2} \right) \ (n \ge 3)$ となる.

また, $D_2(x) = t^2 - t -1$ なので上の式は$n \ge 2$で成り立つ.

したがって, $\alpha := \displaystyle \frac{1 + \sqrt{4n-3}}{2}, \ \beta := \displaystyle \frac{1 - \sqrt{4n-3}}{2}$ とおくと固有値は $0, \ \alpha, \ \beta$.

$P := \begin{pmatrix}
0 & \cdots & 0 & \alpha & \beta \\
-1 & \cdots & -1 & 1 & 1 \\
1 & & & \vdots & \vdots \\
 & \ddots & & \vdots & \vdots \\
 & & 1 & 1 & 1 \end{pmatrix}$
とおくと, 
$P^{-1} = \displaystyle \frac{1}{1-n} 
\begin{pmatrix} 0 & 1 & 2-n & 1 & \cdots & 1 & 1 & 1 \\
                0 & 1 & 1 & 2-n & \cdots & 1 & 1 & 1 \\
 \vdots & \vdots & \vdots & \vdots & \ddots & \vdots & \vdots & \vdots \\
 \vdots & \vdots & \vdots & \vdots & \vdots & \ddots & \vdots & \vdots \\
 \vdots & \vdots & \vdots & \vdots & \vdots & \vdots & \ddots & \vdots \\
               0 & 1 & 1 & 1 & \cdots & 1 & 1 & 2-n \\
\frac{1-n}{\alpha - \beta} & \frac{\beta}{\alpha - \beta} & \cdots & \cdots & \cdots & \cdots & \cdots & \frac{\beta}{\alpha - \beta} \\
\frac{1-n}{\beta - \alpha} & \frac{\alpha}{\beta - \alpha} & \cdots & \cdots & \cdots & \cdots & \cdots & \frac{\alpha}{\beta - \alpha}
\end{pmatrix}$

であり, $P^{-1}AP = \begin{pmatrix}
0 & & & & \\
 & \ddots & & & \\
 & & 0 & & \\
 & & & \alpha & \\
 & & & & \beta \end{pmatrix}$.

\bigskip
\bigskip
\bigskip
\bigskip

(1)
$
\left(\begin{array}{cccc}
2 & 1 & &  \\
1 & \ddots & \ddots &\text{\huge{0}} \\
  & \ddots & \ddots & 1\\
\text{\huge{0}} & & 1 & 2
\end{array}\right)
$
余因子展開


\end{document}
